%*************************************************
% In this file the abstract is typeset.
% Make changes accordingly.
%*************************************************

\addcontentsline{toc}{section}{چکیده}
\newgeometry{left=2.5cm,right=3cm,top=3cm,bottom=2.5cm,includehead=false,headsep=1cm,footnotesep=.5cm}
\setcounter{page}{1}
\pagenumbering{arabic}						% شماره صفحات با عدد
\thispagestyle{empty}

~\vfill

\subsection*{چکیده}
\baselineskip=0.7cm



در دنیای امروز، با پیشرفت فناوری، ابزار‌ها، نیاز‌ه‌ها و فرایند‌ها، روز به روز پیچیده‌تر شده و مدیریت، بهره‌برداری و هزینه مربوط به هرکدام دچار چالش‌های تازه‌ و دشواری می‌شود.
امروزه، مجازی‌سازی اشیاء، موجودیت‌ها و فرایند‌ها، این کمک را می‌کند که بتوان بدون نیاز به ایجاد تغییر و دستکاری در واقعیت، با مشاهده و بررسی دوقلوی دیجیتال آن، بهترین تصمیم گرفته و اجرا شود. 
با این وجود، یک حقیقت غیرقابل انکار این است که یک فرایند واقعی، از ارتباط و تاثیر موجودیت‌ها، اشیاء و فرایند‌های کوچک‌تر بر یکدیگر ایجاد می‌شود و نمی‌توان این مورد را نادیده گرفت. لذا مطالعه بر نوع ارتباط چند عامل، مدیریت،‌ تبادل داده و تصمیم‌گیری درست بر اساس داده‌ها و هوش مصنوعی، می‌تواند به دوقلو‌های دیجیتال در دنیای مجازی کمک کند تا در واقعیت تاثیرگذار باشد.
پس تکنولوژی دوقلو‌ی دجیتال با استفاده از سیستم‌های چند عاملی، می‌تواند توسعه یافته و قابلیت و کارایی بالاتری را در واقعیت فراهم سازد. \\




\begin{small}


\noindent\textbf{کلید ‌واژه‌ها:}
دوقلو‌ی دیجیتال، سیستم‌های چند عاملی، هوش‌ مصنوعی، ابر، مه، لبه، میکروسرویس،
\lr{Eclips Ditto},
\lr{Cintoo}
\\


\end{small}